\documentclass[12pt,a4paper,oneside]{book}
\usepackage[utf8]{inputenc}
\usepackage[T1]{fontenc}
\usepackage{amsmath, amssymb, amsthm, amsfonts}
\usepackage{graphicx}
\usepackage{hyperref}
\usepackage{listings}
\usepackage{xcolor}
\usepackage{geometry}
\usepackage{setspace}
\usepackage{cite}
\usepackage{tikz-cd}

\geometry{
    a4paper,
    left=35mm,
    right=25mm,
    top=25mm,
    bottom=25mm
}

\newtheorem{theorem}{Theorem}[chapter]
\newtheorem{definition}{Definition}[chapter]
\newtheorem{lemma}{Lemma}[chapter]

\definecolor{codegreen}{rgb}{0,0.6,0}
\definecolor{codegray}{rgb}{0.5,0.5,0.5}
\definecolor{codepurple}{rgb}{0.58,0,0.82}
\definecolor{backcolour}{rgb}{0.95,0.95,0.92}

\lstdefinestyle{mystyle}{
    backgroundcolor=\color{backcolour},   
    commentstyle=\color{codegreen},
    keywordstyle=\color{magenta},
    numberstyle=\tiny\color{codegray},
    stringstyle=\color{codepurple},
    basicstyle=\ttfamily\footnotesize,
    breakatwhitespace=false,         
    breaklines=true,                 
    captionpos=b,                    
    keepspaces=true,                 
    numbers=left,                    
    numbersep=5pt,                  
    showspaces=false,                
    showstringspaces=false,
    showtabs=false,                  
    tabsize=2
}
\lstset{style=mystyle}

\begin{document}

\frontmatter

\begin{titlepage}
    \begin{center}
        \vspace*{1cm}
        
        \Huge
        \textbf{The Epistemology of Execution: Functorial Projections from Distributed Observability to Object-Centric Process Reality}
        
        \vspace{1.5cm}
        
        \Large
        \textbf{A Dissertation}
        
        \vspace{1.5cm}
        
        \textit{Submitted in partial fulfillment of the requirements \\ for the degree of Doctor of Philosophy}
        
        \vspace{2cm}
        
        \Large
        Department of Computer Science \\
        Systems, Logic, and Process Intelligence Group
        
        \vspace{1.5cm}
        
        \today
        
    \end{center}
\end{titlepage}

\chapter*{Abstract}
\addcontentsline{toc}{chapter}{Abstract}
The divergence between low-level distributed system observability and high-level object-centric process intelligence represents a fundamental epistemic gap in computer science. Historically, this gap has been bridged by ad-hoc, imperative log parsing, which conflates the observer's instrumentation with the reality of the business process. This thesis establishes a rigorous mathematical and epistemological framework for the deterministic semantic projection of OpenTelemetry (OTEL) execution signals into Object-Centric Event Logs (OCEL). We formalize the observability space $\mathbf{Obs}$ and the process event space $\mathbf{Proc}$ as distinct categories, demonstrating that semantic authority must originate exclusively from public ontologies via declarative morphisms, rather than imperative code. By utilizing graph generators (GGEN) to enforce strict topological boundaries at runtime via Weaver live-checks, we secure the telemetry admission phase. Subsequently, admitted signals are projected via \texttt{SPARQL CONSTRUCT} functors, ensuring that the derivation of process evidence ("what happened") is mathematically insulated from the noise of system telemetry ("what was said"). We prove that this separation yields a fully receiptable, cryptographically verifiable, and logically sound process mining corpus.

\tableofcontents

\mainmatter

\chapter{Introduction}

\section{The Observability-Intelligence Gap}
Modern distributed architectures emit petabytes of telemetry data, primarily structured as logs, metrics, and distributed traces under the OpenTelemetry (OTEL) standard. Concurrently, the field of Process Mining, pioneered by Wil van der Aalst, has evolved from flat, single-case identifier models (like XES - eXtensible Event Stream) to embrace Object-Centric Event Logs (OCEL). OCEL models the complex, many-to-many relationships between process events and physical or logical objects. Flat event logs fail epistemically when mapping microservice telemetry because a single network span often manipulates multiple concurrent domain entities (e.g., a User, an Order, and a Payment). 

However, the translation of raw OTEL spans into highly-structured OCEL corpora relies heavily on imperative parsing scripts. These scripts embed semantic decisions deeply into application code, making the translation fragile, untraceable, and mathematically unprovable.

\section{Thesis Statement}
Imperative code cannot serve as semantic authority. The lawful translation of system observability data into process intelligence requires a deterministic semantic projection pipeline formalized through category theory. By separating telemetry admission from evidence derivation via functorial graph morphisms (\texttt{SPARQL CONSTRUCT}), we can produce mathematically verifiable OCEL evidence directly from OpenTelemetry signals, governed exclusively by public ontologies.

\chapter{Epistemic and Mathematical Foundations}

\section{The Epistemology of Telemetry}
In distributed systems, a trace represents the observer's claim about execution. We define the \textit{Axiom of the Untrusted Observer}: a system's self-reported execution state is an epistemic claim, not an ontological fact. Telemetry represents "what was said." Process intelligence requires "what happened." The transformation between the two must be a formal proof, not a heuristic.

\section{Categorical Formalization and Causal Ordering}
Let $\mathbf{Obs}$ be the category of distributed observations. Drawing upon Lamport's formalization of logical time, the objects in $\mathbf{Obs}$ are OTEL signals (Spans, Events), and the morphisms are strict partial orderings governed by the "happens-before" relation ($\rightarrow$). A distributed trace is therefore a directed acyclic graph (DAG) representing a concurrent execution lattice.

Let $\mathbf{Proc}$ be the category of process reality, where objects are nodes in a bipartite hypergraph consisting of disjoint sets of OCEL Events ($E$) and OCEL Objects ($O$), with edges representing existential and causal qualifications.

\begin{definition}[Semantic Projection Functor]
A Semantic Projection is a structure-preserving functor $F: \mathbf{Obs} \to \mathbf{Proc}$. Because $\mathbf{Obs}$ contains vast operational noise (e.g., latency metrics, CPU utilization) and arbitrary unstructured tags permitted by the Open World Assumption of distributed tracing, $F$ enforces a \textit{Closed World Projection}. It acts as a \textit{forgetful functor} relative to the operational domain, strictly preserving only the ontological invariants defined by the external ontology $\mathcal{O}$, while mapping the happens-before causal lattice of $\mathbf{Obs}$ into the bipartite object-centric causality of $\mathbf{Proc}$. Unmapped variables are rigorously discarded.
\end{definition}

\begin{figure}[h]
\centering
\begin{tikzcd}
    \text{OTEL Span} \arrow[r, "F"] \arrow[d, "\text{Trace Context}"] & \text{OCEL Event} \arrow[d, "\text{ocel:relatedObject}"] \\
    \text{OTEL Attribute} \arrow[r, "F"] & \text{OCEL Object}
\end{tikzcd}
\caption{Commutative Diagram of the Semantic Projection Functor $F$}
\end{figure}

\section{The Axiom of Telemetric Acyclicity}
Distributed clock skew can occasionally violate logical causality, presenting cycles in telemetry. We introduce the \textit{Axiom of Telemetric Acyclicity}: Any sub-graph $g \in \mathbf{Obs}$ must be a strict Directed Acyclic Graph. If $g$ contains a cycle, the morphism $F$ is undefined, and the telemetry must be rejected at the admission boundary prior to graph derivation.

\chapter{Methodology: The Admission Boundary}

\section{The Semantic Compilation Pipeline}
The abstract ontology $\mathcal{O}$ must be made operational without manual interference. We formalize this via the GGEN compiler transformation sequence:
$$ \mathcal{O}_{RDF} \xrightarrow{\text{GGEN}} \Sigma_{Weaver} \xrightarrow{\text{Codegen}} \text{API}_{Rust} $$
Here, $\mathcal{O}_{RDF}$ is the public knowledge graph. GGEN deterministically translates this into a topological bounding surface $\Sigma_{Weaver}$ (the Weaver semantic-convention registry). Finally, Weaver's code generator produces strongly-typed APIs, structurally restricting the observer from emitting undefined propositions.

\section{Topological Boundaries via Weaver}
To prevent semantic corruption in $\mathbf{Obs}$, we establish a strict topological boundary. GGEN (Graph Generator) projects $\mathcal{O}$ into a Weaver semantic-convention registry. The Weaver \texttt{live-check} acts as the boundary enforcement mechanism.

\begin{theorem}[Admission Completeness via Entailment]
Let $\Sigma_t$ be the set of semantic constraints projected from $\mathcal{O}_t$ into the Weaver registry at schema version $t$. An observation sub-graph $g \in \mathbf{Obs}$ is admitted if and only if $g \models \Sigma_t$, where $\models$ denotes semantic entailment under the Weaver topological boundary. Consequently, the state space of admitted telemetry is a strict topological subspace of mathematically valid propositions, immune to runtime structural poisoning.
\end{theorem}

\section{Ontological Polymorphism and Schema Evolution}
Real distributed systems undergo constant schema evolution. When $\mathcal{O}_t$ evolves to $\mathcal{O}_{t+1}$, we define a schema morphism $S: \mathcal{O}_t \to \mathcal{O}_{t+1}$. The projection functor $F$ is polymorphic over $S$, guaranteeing that legacy observations $g_t$ can be deterministically projected into the current process reality hypergraph provided $S$ forms a valid natural transformation.

\chapter{Functorial Materialization: OTEL to OCEL}

\section{SPARQL as the Morphism Engine}
The functor $F$ is materialized not through imperative code, but through \texttt{SPARQL CONSTRUCT} over an RDF representation of the admitted $G_{OTEL}$. 

\begin{lstlisting}[language=SQL, caption=Morphism $F$ expressed in SPARQL]
CONSTRUCT {
  ?event a ocel:Event ;
         ocel:eventType ?eventType ;
         ocel:relatedObject ?execution .
}
WHERE {
  ?span a otel:Span ;
        otel:traceId ?traceId ;
        otel:spanId ?spanId ;
        otel:attribute [ otel:key "process.activity.iri" ; otel:value ?activity ] ;
        otel:attribute [ otel:key "process.object.id" ; otel:value ?objectId ] .
  
  BIND(IRI(CONCAT("urn:ocel:event:", STR(?spanId))) AS ?event)
  BIND(IRI(CONCAT("urn:ocel:object:", STR(?objectId))) AS ?execution)
}
\end{lstlisting}

\section{Hypergraph Edge Construction and Refusals}
Crucially, the morphism $F$ does not merely instantiate entities; it constructs the bipartite edges of $\mathbf{Proc}$. The `ocel:relatedObject` predicate acts as the incidence function mapping events to objects. In the example above, the OTEL attribute `process.object.id` dynamically binds the process execution event to its participating object, forming the hypergraph topology without imperative heuristics.

Furthermore, process intelligence must model negative executions. If an observation represents an aborted transaction (e.g., an OTEL Span with \texttt{StatusCode::Error}), the functor $F$ materializes a \textit{Terminal State Refusal} edge in $\mathbf{Proc}$, representing an ontologically valid halted workflow, distinct from a silent system failure.

\section{Computational Complexity of Materialization}
A rigorous projection must scale to billions of observational events. Because $F$ operates over a DAG of admitted traces, the evaluation of the SPARQL \texttt{CONSTRUCT} queries over the RDF graph $G_{OTEL}$ can be bounded. Assuming a topologically sorted trace graph, the worst-case time complexity of materializing $G_{OCEL}$ is strictly bounded by $O(|V_{OTEL}| \log |V_{OTEL}| + |E_{OTEL}|)$, guaranteeing tractability for extreme-scale enterprise telemetry processing.


\chapter{Cryptographic Provenance and Idempotent Replay}

\section{SHACL Validation and Conformance}
Once $F$ materializes the object-centric evidence $G_{OCEL}$, the graph must be formally validated against the structural constraints of the business process. Let $\mathcal{S}$ represent the Shapes Constraint Language (SHACL) definitions derived from $\mathcal{O}$. Process conformance is mathematically defined as the evaluation of $G_{OCEL}$ against $\mathcal{S}$. If $G_{OCEL} \models \mathcal{S}$, the process execution is deemed conformant.

\section{Idempotent Process Replay}
A defining property of the semantic projection is its reproducibility. Because $G_{OTEL}$ is an immutable, receipted observation graph, and $P$ is a pure functional morphism, the replay operation is strictly idempotent. Re-evaluating $P(G_{OTEL})$ guarantees an exact graph isomorphism: $G_{OCEL} \cong G'_{OCEL}$ for all subsequent executions. This eliminates the non-determinism inherent in parsing transient, mutable log files.

\section{The Non-Repudiable Projection}
Because the projection $F$ is completely deterministic and relies solely on declarative morphisms, the derivation of $G_{OCEL}$ is strictly receiptable. Let $H(\cdot)$ denote a collision-resistant cryptographic hash function. The materialization of process intelligence is defined as a verifiable provenance chain:

\begin{definition}[Cryptographic Provenance Chain]
The process reality graph $G_{OCEL}$ is bound to its origin through the hash commitment $H(G_{OCEL})$, satisfying the invariant:
$$ H(G_{OCEL}) = H \Big( H(G_{OTEL}) \parallel H(P) \parallel H(\mathcal{O}) \Big) $$
where $P$ is the SPARQL morphism artifact and $\mathcal{O}$ is the public ontology.
\end{definition}

This guarantees that any tampering with the runtime telemetry ($G_{OTEL}$), the projection logic ($P$), or the ontological definitions ($\mathcal{O}$) invalidates the signature of the process reality, yielding absolute non-repudiation in process mining audits.

\chapter{Conclusion: The Grand Unification}
This dissertation formalizes a mathematical pipeline connecting the lowest levels of distributed systems observability with the highest levels of object-centric process intelligence. We propose the \textit{Grand Unification Equation of Process Observability}:
$$ \mathbf{Proc} = \text{SHACL\_VALIDATE} \Big( F_{SPARQL} \big(\text{WEAVER\_ADMIT}( \mathbf{Obs}, \mathcal{O} ) \big) \Big) $$

By treating ontologies as the sole semantic authority ($\mathcal{O}$), validating telemetry via Weaver topological bounds, and utilizing declarative graph projections ($F_{SPARQL}$) instead of imperative parsers, we establish a verifiable, receiptable, and completely deterministic path from OpenTelemetry to OCEL. Empirical implementation of this mathematics was achieved utilizing the \texttt{chicago-tdd-tools} telemetry admission gateways and the \texttt{wasm4pm} Oxigraph semantic reasoning engine, definitively closing the epistemic gap in process mining.

\backmatter
\begin{thebibliography}{9}
\bibitem{weber2023}
Weber, I., et al. \textit{Object-Centric Process Mining}. Springer, 2023.
\end{thebibliography}

\end{document}
